\section{Week 3 - 28/9}
\subsection{Resummation for current expression}
Firstly, we'd like to find a general expression for the current in the model, without approximating the noise. For this we solve the self consistent set of equations given by \autoref{eq:c-evolution} and \autoref{eq:d-evolution}. Transforming to Fourier space we find
\begin{eqnarray}\label{eq:c-fourier}
    \hat{c}_k(E) = \hat{c}_{k,0}(E) + T G_{k,0}^{(+)}(E)\hat{d}(E), \\ \label{eq:d-fourier}
    \hat{d}(E) = \hat{d}_0(E) + T G_{S,0}^{(+)}(E)\sum_k \hat{c}_k(E).
\end{eqnarray}
Focussing on the equation for $\hat{d}(E)$, we substitute \ref{eq:c-fourier}, giving 
\begin{equation}
    \hat{d}(E) = \hat{d}_0(E) + T G_{S,0}^{(+)}(E)\sum_k \hat{c}_{k,0}(E) + T^2 G_{S,0}^{(+)}(E)\sum_k G_{k,0}^{(+)}(E)\hat{d}(E),
\end{equation}
which is solved by
\begin{equation}
    \hat{d}(E) = \frac{\hat{d}_0(E) + T G_{S,0}^{(+)}(E)\sum_k \hat{c}_{k,0}(E)}{1- T^2 G_{S,0}^{(+)}(E)\sum_k G_{k,0}^{(+)}(E)} = \left(1+\hbar\Sigma_S^{(+)}(E)\right)\hat{d}_0(E)+TG_{S}^{(+)}(E)\sum_k\hat{c}_{k,0}(E),
\end{equation}
where $G_S^{(+)}(E) = ((G_{S,0}^{(+)})^{-1} - T^2\sum_kG_{k,0}^{(+)}(E))^{-1}$ is the total retarded Green's function in the sample, and the retarded self-energy of the sample is given by $\hbar\Sigma^{(+)}_S(E) = T^2\sum_k G_{k,0}^{(+)}(E)$. Similarly, for $\hat{c}_k(E)$ we can substitute \autoref{eq:d-fourier} in \autoref{eq:c-fourier}, so that
\begin{equation}
    \hat{c}_k(E) = \hat{c}_{k,0}(E) + T G_{k,0}^{(+)}(E) \hat{d}_0(E) + T^2 G_{k,0}^{(+)}(E)G_{S,0}^{(+)}(E)\sum_q \hat{c}_q(E),
\end{equation}
which is less trivial to solve. Our system doesn't conserve momentum, due to a lack of translational invariance. It does conserve energy, as the Hamiltonian is symmetric under time translations. The absence of momentum conservation means that electrons of different momenta may couple when they tunnel back and forth through the sample. This is also what you can see in the above equation. To solve the equation we write this as 
\begin{equation}
    \sum_q \left[\delta_{kq} - T^2G_{k,0}^{(+)}(E)G_{S,0}^{(+)}(E)\right]\hat{c}_q(E) = \hat{c}_{k,0}(E) + T G_{k,0}^{(+)}(E) \hat{d}_0(E).
\end{equation}
Which is inverted by using $\sum_l[\delta_{kl} + T^2 G_{k,0}^{(+)}(E)G_S^{(+)}(E)][\delta_{lq} - T^2G_{l,0}^{(+)}(E)G_{S,0}^{(+)}(E)] = \delta_{kq}$ as is readily checked (note the total Green's function)
\newline
\begin{align}
    \hat{c}_k(E) &= \sum_l \left[\delta_{kl} + T^2 G_{k,0}^{(+)}(E)G_S^{(+)}(E)\right] \left(\hat{c}_{l,0}(E) + T G_{l,0}^{(+)}(E) \hat{d}_0(E)\right), \nonumber\\
    &= \hat{c}_{k,0}(E) + T G_{k,0}^{(+)}(E) \hat{d}_0(E) + T^2 G_{k,0}^{(+)}(E)G_S^{(+)}(E)\sum_q\hat{c}_{q,0}(E) \nonumber\\
    & \quad + T^3 G_{k,0}^{(+)}(E)G_S^{(+)}(E)\sum_qG_{q,0}^{(+)}(E)\hat{d}_0(E).
\end{align}
Combining the second and fourth term we get
\begin{align}
    \hat{c}_k(E) &= \hat{c}_{k,0}(E) + T^2 G_{k,0}^{(+)}(E)G_S^{(+)}(E)\sum_q\hat{c}_{q,0}(E) + TG_{k,0}^{(+)}(E) \underbrace{\left(1 + T^2 G_S^{(+)}(E)\sum_q G_{q,0}^{(+)}(E)\right)}_{=\left(1+\hbar\Sigma^{(+)}_S(E)G_S^{(+)}(E)\right)}\hat{d}_0(E),\nonumber\\
    &= \hat{c}_{k,0}(E) + T^2 G_{k,0}^{(+)}(E)G_S^{(+)}(E)\sum_q\hat{c}_{q,0}(E)+ TG_{k,0}^{(+)}(E)\left(1+\hbar\Sigma^{(+)}_S(E)G_S^{(+)}(E)\right)\hat{d}_0(E).
\end{align}
We can now fill these equations in into $\langle \hat{d}^\dagger(t)\hat{c}_k(t)\rangle$, 
\begin{equation}
    \langle \hat{d}^\dagger(t)\hat{c}_k(t)\rangle = \int\frac{dE}{2\pi\hbar}\frac{dE'}{2\pi\hbar}e^{-\frac{i(E'-E)t}{\hbar}}\langle\hat{d}^\dagger(E)\hat{c}_k(E')\rangle.
\end{equation}
The Hamiltonian is has time translation symmetry, so that the correlators are energy-conserving, i.e. $\langle \hat{d}^\dagger(E)\hat{c}_k(E')\rangle = 2\pi\hbar\delta(E-E')\langle \hat{d}^\dagger(E)\hat{c}_k(E)\rangle$. The crossterm correlators of the free operators are zero by definition, including the crossterms of different tip momenta, i.e. for $k\neq q$ we have $\langle \hat{c}_{k,0}(E)\hat{c}_{q,0}(E)\rangle =0$. This means we get
\begin{align}
    \langle \hat{d}^\dagger(t)\hat{c}_k(t)\rangle &= \int\frac{dE}{2\pi\hbar}\langle\hat{d}^\dagger(E)\hat{c}_k(E)\rangle,\nonumber\\
    &=\int\frac{dE}{2\pi\hbar}\left[\left(1 + \hbar\Sigma_S^{(-)}(E)G_S^{(-)}(E)\right)TG_{k,0}^{(+)}(E)\left(1+\hbar\Sigma^{(+)}_S(E)G_S^{(+)}(E)\right)\langle\hat{d}^\dagger_0(E)\hat{d}_0(E)\rangle\right. \\
    &\left.+ TG_{S}^{(-)}(E)\langle\hat{c}^\dagger_{k,0}(E)\hat{c}_{k,0}(E)\rangle+T^3G_{S}^{(-)}(E)G_{k,0}^{(+)}(E)G_S^{(+)}(E)\sum_q\langle\hat{c}^\dagger_{q,0}(E)\hat{c}_{q,0}(E)\rangle\right]
\end{align}








\subsection{Symmetrizing the noise}
The one-sided power spectrum is properly defined as
\begin{equation}
    \tilde{S}(E) = \int d\tau \left(S(\tau) + S(-\tau)\right)e^{\frac{iE\tau}{\hbar}} = \int d\tau\,S(\tau)\left(e^{\frac{iE\tau}{\hbar}}+e^{-\frac{iE\tau}{\hbar}}\right),
\end{equation}
i.e. the noise should be symmetrized. This gives
\begin{align}
    \tilde{S}(E) &= e^2|T|^2 \sum_k \int d\tau \int \frac{dE'}{2\pi\hbar} \int\frac{dE''}{2\pi\hbar} \left[G^>_{k,0}(E')G^<_{S,0}(E'')e^{-\frac{i\tau(E'-E'')}{\hbar}} \right. \nonumber \\ 
    &\left.\quad+ G^<_{k,0}(E')G^>_{S,0}(E'')e^{-\frac{i\tau(E''-E')}{\hbar}}\right]\left(e^{\frac{iE\tau}{\hbar}}+e^{-\frac{iE\tau}{\hbar}}\right) + O(T^3), \nonumber\\
    &=e^2|T|^2 \sum_k \int \frac{dE'}{2\pi\hbar}\left(G^>_{k,0}(E')\left[G^<_{S,0}(E'-E) + G^<_{S,0}(E'+E)\right]+\right.\nonumber\\
    &\left. \quad G^<_{k,0}(E')\left[G^>_{S,0}(E'-E) + G^>_{S,0}(E'+E)\right]\right)+ O(T^3).
\end{align}
Using fluctuation-dissipation we get
\begin{align}
    \tilde{S}(E) = \frac{2\pi e^2|T|^2}{\hbar} \int dE'\,\rho_N(E')\biggl((1-N_N(E'))\bigl[\rho_S(E'-E)N_S(E'-E) + \rho_S(E'+E)N_S(E'+E)\bigr]\nonumber\\
    + N_N(E')\bigl[\rho_S(E'-E)(1-N_S(E'-E)) + \rho_S(E'+E)(1-N_S(E'+E))\bigr]\biggr) + O(T^3).
\end{align}
We introduce the tunneling rate $\Gamma_T(E) = \frac{2\pi}{\hbar}|T|^2 \rho_N(E) \approx \Gamma_T$, which we take to be constant, as the density of states of the normal conductor is slowly varying around the Fermi energy. At zero temperature the distribution function reduces to a Heaviside step function centered at the Fermi energy, but in the normal metal the Fermi energy is shifted by the applied voltage, so that the step lies at $eV/2$, giving
\begin{equation}
    N_N(E) = \theta\left(\frac{eV}{2}-E\right); \qquad N_S(E) = \theta\left(-E\right).
\end{equation}
The factor $1/2$ is because we're using a one-lead model. Using $1-\theta(x) = \theta(-x)$ we get
\begin{align}
    \tilde{S}(E) = e^2\Gamma_T \int dE'\,\left(\theta\left(E'-\frac{eV}{2}\right)\left[\rho_S(E'-E)\theta\left(E-E'\right) + \rho_S(E'+E)\theta\left(-E-E'\right)\right]\right.\nonumber\\
    \left.+ \theta\left(\frac{eV}{2}-E'\right)\left[\rho_S(E'-E)\theta\left(E'-E\right) + \rho_S(E'+E)\theta\left(E'+E\right)\right]\right) + O(T^3).
\end{align}
We recognize each term has a lower and upper bound on $E'$, so we can simplify further to 
\begin{align}
    \tilde{S}(E) &=e^2 \Gamma_T \left(\theta\left(E-\frac{eV}{2}\right)\int_{-(E-eV/2)}^0dE'\,\rho_S(E') + \theta\left(-E-\frac{eV}{2}\right)\int_{-(-E-eV/2)}^0dE'\,\rho_S(E')\right.\nonumber\\ 
    &\left.\qquad\qquad\theta\left(\frac{eV}{2}-E\right)\int_0^{eV/2-E}dE'\,\rho_S(E') + \theta\left(\frac{eV}{2}+E\right)\int_0^{eV/2+E}dE'\,\rho_S(E')\right) + O(T^3).
\end{align}
Which simplifies further to (because the density of states is strictly positive)
\begin{align}
    \tilde{S}(E) &=e^2 \Gamma_T \left(\left|\int_0^{eV/2-E}dE'\,\rho_S(E')\right|+ \left|\int_0^{eV/2+E}dE'\,\rho_S(E')\right|\right)+O(T^3).
\end{align}
Where at this order we just have the free density of states. However, when we replace the free Green's function by the total Green's functions $G^{(+)}(E)$ we find the full density of states, which is a Lorentzian, given by 
\begin{equation}
    \rho_S(E) = \frac{i}{2\pi}\left[\frac{1}{E-\xi_S-\hbar\Sigma} - \frac{1}{E-\xi_S-\hbar\Sigma^*}\right] = \frac{1}{\pi}\frac{-\hbar\mathfrak{Im}(\Sigma)}{(E-\xi_S-\hbar\mathfrak{Re}(\Sigma))^2+\hbar^2\mathfrak{Im}^2(\Sigma)},
\end{equation}
where $\Sigma$ is the self-energy of the sample, given by
\begin{equation}
    \Sigma = |T|^2 \sum_k\frac{1}{E-\xi_k+i\eta}.
\end{equation}
The real and imaginary parts are then given by (after $\eta \downarrow 0$)
\begin{gather}
    \hbar\mathfrak{Re}(\Sigma) = |T|^2\sum_k\frac{1}{E-\xi_k},\\
    \hbar\mathfrak{Im}(\Sigma) = -\pi|T|^2\sum_k\delta(E-\xi_k) = -\pi|T|^2\rho_N(E) = -\frac{\hbar\Gamma_T}{2}.
\end{gather}
We can ignore the shift to the energy of the state due to the real part, as it's small \textcolor{red}{(although I'm not sure why)}. The imaginary part we can relate to the tunneling rate $\Gamma_T$ defined prior, which we take to be energy independent. We find the density of states given by
\begin{equation}
    \rho_S(E) = \frac{\hbar\Gamma_T}{2\pi}\frac{1}{(E-\xi_S)^2+(\hbar\Gamma_T/2)^2} = \frac{2}{\pi\hbar\Gamma_T}\frac{1}{\left(\frac{2(E-\xi_S)}{\hbar\Gamma_T}\right)^2+1}.
\end{equation}
Which we can explicitly integrate
\begin{align}
    \tilde{S}(E) &= \frac{2e^2}{\pi\hbar}\left(\left|\int_0^{eV/2-E}dE'\,\frac{1}{\left(\frac{2(E'-\xi_S)}{\hbar\Gamma_T}\right)^2+1}\right|+\left|\int_0^{eV/2+E}dE'\frac{1}{\left(\frac{2(E'-\xi_S)}{\hbar\Gamma_T}\right)^2+1}\right|\right),\\
    &=\frac{e^2\Gamma_T}{\pi}\left(\left|\arctan\left[\frac{2}{\hbar \Gamma_T}(eV/2-E-\xi_S)\right]-\arctan\left[-\frac{2}{\hbar \Gamma_T}\xi_S\right]\right|\right.\\
    &\left.\qquad\qquad+\left|\arctan\left[\frac{2}{\hbar \Gamma_T}(eV/2+E-\xi_S)\right]-\arctan\left[-\frac{2}{\hbar \Gamma_T}\xi_S\right]\right|
    \right).
\end{align}
The current is given by \autoref{eq:current}, which at zero temperature and the above approximations reduces to
\begin{align}
    \langle \hat{I} \rangle &= e\Gamma_T\int_0^{eV/2}dE\,\rho_S(E),\\
    &=\frac{2e}{\pi\hbar}\int_0^{eV/2}dE\,\frac{1}{\left(\frac{2(E-\xi_S)}{\hbar\Gamma_T}\right)^2+1},\\
    &=\frac{e\Gamma_T}{\pi}\left(\arctan\left[\frac{2}{\hbar \Gamma_T}(eV/2-\xi_S)\right]-\arctan\left[-\frac{2}{\hbar \Gamma_T}\xi_S\right]\right).
\end{align}
Using this we can plot the Fano factor as a function of energy (at zero temperature there is no thermal noise). 
\begin{tcolorbox}[title = Questions/Ideas, colback = gray!5, colframe=black]
    \begin{enumerate}
        \item The graph of the noise spectrum dips below its value at zero, which means that the shot noise can't be white noise. Why not?
        \item I'm not exactly sure what the results would be if we would explicitly carry out the resummation, but I'm not sure you'd just get the total Green's function.
        \item \textcolor{blue}{Use $G_0^{-1} = G^{-1}+\hbar\Sigma$ to further simplify the expression for $\langle\hat{d}^\dagger\hat{c}_k\rangle$}
        \item \textcolor{blue}{Calculate the Fano factor for large energy in the Lorentzian case.}
        \item \textcolor{blue}{Make graphs of the noise spectrum for constant density, and compare to paper by Schoelkopf.}
        \item \textcolor{blue}{Do calculation of noise spectrum for a superconducting DOS.}
        \item \textcolor{blue}{Do a fourth order calculation, or a time space calculation. What is the groundstate, and how does the evolution work out (do we take $t_0 \rightarrow -\infty$)?}
    \end{enumerate}
\end{tcolorbox}

\begin{tcolorbox}[title = Sources, colback = gray!5, colframe=gray]
    \begin{enumerate}
        \item Martins - Noise in mesoscopic physics
        \item Schoelkopf - Frequency Dependence of Shot Noise in a Diffusive Mesoscopic Conductor
    \end{enumerate}
\end{tcolorbox}

\begin{tcolorbox}[title = Roadmap, colback = red!7!white, colbacktitle = red!95!black]
    There's two separate calculations going on. The first is the resummation of the Green's function, the second is the calculation of the noise spectrum. 
\end{tcolorbox}

